% =============================================================================
%                    WEEK 1: NON-COOPERATIVE GAMES (PART I)
% =============================================================================
\section{Week 1: Non-Cooperative Games (Part I)}

% -----------------------------------------------------------------------------
%                              1.1 TOPIC SUMMARY
% -----------------------------------------------------------------------------
\subsection{Topic Summary}

\begin{keybox}[Learning Objectives]
By the end of this week, you will be able to:
\begin{enumerate}
  \item Define the core concepts of non-cooperative game theory, including \textbf{strategy profiles}, \textbf{payoff functions}, and the distinction between cooperative and non-cooperative games.
  \item Formulate real-world scenarios as \textbf{strategic-form (normal-form)} or extensive-form games involving rational, self-interested agents.
  \item Identify and classify various types of strategies: \textbf{pure}, \textbf{mixed}, \textbf{dominant}, and \textbf{best response} strategies.
  \item Analyse game-theoretic models to determine the existence and properties of \textbf{Nash equilibria} (in pure and mixed strategies).
  \item Explain and interpret equilibrium behaviour in classical games: \textbf{Prisoner's Dilemma}, \textbf{Matching Pennies}, \textbf{Battle of the Sexes}, \textbf{Stag Hunt}, and coordination games.
  \item Evaluate the potential for and limitations of non-cooperative solutions, including issues of \textbf{uniqueness}, \textbf{multiplicity}, and \textbf{inefficiency} of equilibria.
  \item Demonstrate competency in using mathematical notation and reasoning to represent and solve strategic games.
\end{enumerate}
\end{keybox}

\separator

\subsubsection*{The Big Picture}

This week introduces \textbf{non-cooperative game theory}---the mathematical framework for analysing strategic interactions between rational decision-makers. The module sits at the intersection of two fields:

\begin{itemize}
  \item \textbf{Algorithmic Game Theory (AGT):} Studies algorithms in strategic environments where inputs come from self-interested agents. Combines theoretical CS and economics to analyse equilibria and design efficient, fair mechanisms.
  \item \textbf{Computational Social Choice (COMSOC):} Studies collective decision-making and preference aggregation, asking how individual preferences can be combined into collective decisions.
\end{itemize}

The central question: \textit{``What is a stable outcome when multiple self-interested agents interact?''}

\separator

\subsubsection*{What Examiners Like to Ask}

\begin{itemize}
  \item Define a game in \textbf{normal form}; write the formal tuple $\Gamma = (N, (S_i), (U_i))$.
  \item Given a payoff matrix, identify \textbf{dominated strategies} and solve by iterated elimination.
  \item Find all \textbf{pure strategy Nash equilibria} using best-response analysis.
  \item Compute \textbf{mixed strategy Nash equilibria} using the indifference principle.
  \item Compare \textbf{maximin} vs \textbf{minimax} strategies; apply the Minimax theorem to zero-sum games.
  \item Explain why Prisoner's Dilemma has a unique NE that is Pareto-inefficient.
  \item Distinguish \textbf{Nash equilibrium} from \textbf{strong equilibrium}.
\end{itemize}

\separator

% -----------------------------------------------------------------------------
%                 1.2 OVERVIEW + CORE CONCEPTS + EXTENDED THEORY
% -----------------------------------------------------------------------------
\subsection{Overview + Core Concepts + Extended Theory}

\subsubsection*{Roadmap}
\begin{enumerate}
  \item Multi-agent systems and game types
  \item Normal form games: formal definition
  \item Dominance and iterated elimination
  \item Nash equilibrium in pure strategies
  \item Mixed strategies and expected payoffs
  \item Nash equilibrium in mixed strategies
  \item Computing mixed NE (indifference principle)
  \item Maximin and minimax strategies
  \item Zero-sum games and the Minimax theorem
  \item Other solution concepts (strong equilibrium, correlated equilibrium)
\end{enumerate}

\bigseparator

% --- Multi-Agent Systems ---
\subsubsection{Multi-Agent Systems}

\begin{defbox}[Multi-Agent System (MAS)]
A \textbf{multi-agent system} consists of multiple autonomous entities (agents) with:
\begin{itemize}
  \item Distributed information
  \item Computational capabilities
  \item Potentially conflicting objectives
\end{itemize}

\textbf{Types:}
\begin{center}
\renewcommand{\arraystretch}{1.3}
\begin{tabular}{@{} l p{8cm} @{}}
\toprule
\textbf{Type} & \textbf{Examples} \\
\midrule
Cooperative & Sensor networks, robot teams in warehouses \\
Competitive & Electronic marketplaces, resource allocation \\
Mixed & Wikipedia, Stack Overflow \\
\bottomrule
\end{tabular}
\end{center}
\end{defbox}

\separator

% --- Normal Form Games ---
\subsubsection{Normal Form Games}

\begin{defbox}[Non-Cooperative Game]
In a \textbf{non-cooperative} (or competitive) game:
\begin{itemize}
  \item Players take actions and enjoy outcomes \textit{individually}
  \item There are \textbf{no binding contracts} among players
  \item Each player strategically optimises their own utility function
\end{itemize}
\end{defbox}

\begin{defbox}[Normal Form Game (Formal Definition)]
A game in \textbf{strategic (normal) form} is a tuple:
\[
\Gamma = (N, (S_i)_{i \in N}, (U_i)_{i \in N})
\]

where:
\begin{itemize}
  \item $N$ is a finite set of \textbf{players}
  \item $S_i$ is the set of \textbf{strategies} for player $i \in N$
  \item $U_i : S \to \mathbb{R}$ is the \textbf{utility function} for player $i$, where $S = \times_{i \in N} S_i$ is the set of all strategy profiles
\end{itemize}
\end{defbox}

\begin{defbox}[Strategy Profile Notation]
A \textbf{profile} $s$ is a vector of strategies $(s_1, \ldots, s_n)$ where $s_i \in S_i$.

\textbf{Key notation:}
\begin{itemize}
  \item $s_{-i} = (s_1, \ldots, s_{i-1}, s_{i+1}, \ldots, s_n)$ --- profile without player $i$'s strategy
  \item $s = (s_i, s_{-i})$ --- decomposition of profile
  \item $s_C$ --- strategies of coalition $C \subseteq N$
  \item $s_{-C}$ --- strategies of players not in $C$
\end{itemize}
\end{defbox}

\begin{tipbox}[Exam Tip: Writing Game Definitions]
When asked to ``formulate as a game'', always specify:
\begin{enumerate}
  \item The set of players $N$
  \item The strategy set $S_i$ for each player
  \item The utility function $U_i$ (or payoff matrix for finite games)
\end{enumerate}
\end{tipbox}

\separator

% --- Prisoner's Dilemma ---
\subsubsection{The Prisoner's Dilemma}

\begin{exbox}[Prisoner's Dilemma]
\textbf{Setup:} Two suspects are interrogated in separate rooms. Each can:
\begin{itemize}
  \item \textbf{Cooperate (C):} Deny the crime
  \item \textbf{Defect (D):} Confess (rat out the other)
\end{itemize}

\textbf{Payoff matrix:}
\begin{center}
\renewcommand{\arraystretch}{1.3}
\begin{tabular}{@{} c | c c @{}}
 & C & D \\
\midrule
C & $-1, -1$ & $-10, 0$ \\
D & $0, -10$ & $-5, -5$ \\
\end{tabular}
\end{center}

\textbf{Reading the matrix:} First number = row player's payoff; second = column player's.

\textbf{Key insight:} The only stable outcome is $(D, D)$, even though $(C, C)$ gives both players a better payoff!
\end{exbox}

\begin{defbox}[Prisoner's Dilemma: General Form]
\begin{center}
\renewcommand{\arraystretch}{1.3}
\begin{tabular}{@{} c | c c @{}}
 & C & D \\
\midrule
C & $R, R$ & $S, T$ \\
D & $T, S$ & $P, P$ \\
\end{tabular}
\end{center}

\textbf{Condition:} $T > R > P > S$

\begin{itemize}
  \item $T$ = Temptation (defect while other cooperates)
  \item $R$ = Reward (mutual cooperation)
  \item $P$ = Punishment (mutual defection)
  \item $S$ = Sucker's payoff (cooperate while other defects)
\end{itemize}
\end{defbox}

\bigseparator

% --- Dominance ---
\subsubsection{Dominance}

\begin{defbox}[Strict and Weak Dominance]
For player $i$ and strategies $s, t \in S_i$:

\textbf{Strict dominance:} $s$ \term{strictly dominates} $t$ if:
\[
\forall s_{-i} \in S_{-i}: \quad U_i(s, s_{-i}) > U_i(t, s_{-i})
\]

\textbf{Weak dominance:} $s$ \term{weakly dominates} $t$ if:
\[
\forall s_{-i} \in S_{-i}: U_i(s, s_{-i}) \geq U_i(t, s_{-i}) \quad \text{and} \quad \exists s'_{-i}: U_i(s, s'_{-i}) > U_i(t, s'_{-i})
\]
\end{defbox}

\begin{defbox}[Dominant and Dominated Strategies]
\begin{itemize}
  \item Strategy $s$ is \textbf{dominant} if it dominates \textit{all} other strategies
  \item Strategy $t$ is \textbf{dominated} if \textit{some} other strategy dominates it
\end{itemize}

\textbf{Key principle:} Rational players never choose strictly dominated strategies.
\end{defbox}

\begin{tipbox}[Exam Tip: Dominance is Rare but Powerful]
Dominant strategy is a very strong solution concept---no assumptions needed about other players. However, it rarely exists. In Prisoner's Dilemma, Defect is dominant for both players.
\end{tipbox}

\separator

% --- Iterated Elimination ---
\subsubsection{Iterated Elimination of Dominated Strategies (IEDS)}

\begin{thmbox}[IEDS Algorithm]
\textbf{Procedure:}
\begin{enumerate}
  \item Find a strictly dominated strategy for any player
  \item Eliminate it from the game
  \item Repeat until no dominated strategies remain
\end{enumerate}

\textbf{Properties:}
\begin{itemize}
  \item \textbf{Strict dominance:} Order-independent (same final result)
  \item \textbf{Weak dominance:} Order-dependent (different elimination orders may give different results)
  \item Surviving strategies need not be dominant
  \item May end with multiple strategies per player
\end{itemize}
\end{thmbox}

\begin{exbox}[Example: IEDS on a $3 \times 3$ Game]
\textbf{Original game:}
\begin{center}
\renewcommand{\arraystretch}{1.2}
\begin{tabular}{@{} c | c c c @{}}
 & L & C & R \\
\midrule
T & $73, 25$ & $57, 42$ & $66, 32$ \\
M & $80, 26$ & $35, 12$ & $32, 54$ \\
B & $28, 27$ & $63, 31$ & $54, 29$ \\
\end{tabular}
\end{center}

\textbf{Step 1:} For column player, R dominates L (compare: $32 > 25$, $54 > 26$, $29 > 27$). Eliminate L.

\textbf{Step 2:} For row player, T and B both dominate M. Eliminate M.

\textbf{Step 3:} For column player, C dominates R. Eliminate R.

\textbf{Step 4:} For row player, B dominates T ($63 > 57$). Eliminate T.

\textbf{Solution:} $(B, C)$ with payoffs $(63, 31)$.

\textbf{Verify NE:} Row player: $63 > 54$ (B vs T against C). Column player: $31 > 29$ (C vs R against B). \checkmark
\end{exbox}

\bigseparator

% --- Nash Equilibrium ---
\subsubsection{Nash Equilibrium (Pure Strategies)}

\begin{defbox}[Nash Equilibrium]
A strategy profile $s^* = (s^*_1, \ldots, s^*_n)$ is a \textbf{(pure strategy) Nash equilibrium} if no player has a \textbf{unilateral improving move}:
\[
\forall i \in N, \forall s'_i \in S_i: \quad U_i(s^*_i, s^*_{-i}) \geq U_i(s'_i, s^*_{-i})
\]

Equivalently: each player plays a \textbf{best response} to the others' strategies.
\end{defbox}

\begin{defbox}[Best Response]
Strategy $s_i$ is a \textbf{best response} to $s_{-i}$ if:
\[
s_i \in \arg\max_{s'_i \in S_i} U_i(s'_i, s_{-i})
\]
\end{defbox}

\begin{tipbox}[Exam Tip: Finding Pure Strategy NE]
\textbf{Method 1 (Best Response):}
\begin{enumerate}
  \item For each opponent strategy, find player's best response (underline/mark it)
  \item NE = cells where \textit{all} players are playing best responses
\end{enumerate}

\textbf{Method 2 (Check deviations):}
\begin{enumerate}
  \item For each cell, check if any player wants to deviate unilaterally
  \item NE = cells with no profitable deviations
\end{enumerate}
\end{tipbox}

\begin{exbox}[Example: Best Response Analysis]
\begin{center}
\renewcommand{\arraystretch}{1.2}
\begin{tabular}{@{} c | c c c @{}}
 & L & C & R \\
\midrule
T & $2, 2$ & $1^*, 3^*$ & $0^*, 1$ \\
M & $3^*, 1^*$ & $0, 0$ & $0^*, 0$ \\
B & $1, 0^*$ & $0, 0^*$ & $0^*, 0^*$ \\
\end{tabular}
\end{center}

\textbf{Reading:} $^*$ marks best responses. NE = cells where both payoffs are marked.

\textbf{Nash equilibria:} None! No cell has both payoffs marked.
\end{exbox}

\separator

% --- Coordination Games ---
\subsubsection{Coordination Games}

\begin{defbox}[Coordination Game]
A game where players benefit from \textbf{choosing corresponding strategies}:
\begin{center}
\renewcommand{\arraystretch}{1.2}
\begin{tabular}{@{} c | c c @{}}
 & Left & Right \\
\midrule
Up & $1, 1$ & $0, 0$ \\
Down & $0, 0$ & $1, 1$ \\
\end{tabular}
\end{center}

\textbf{Properties:}
\begin{itemize}
  \item \textbf{Two pure strategy NE:} (Up, Left) and (Down, Right)
  \item No dominant or dominated strategies
  \item Both NE are also \textbf{strong equilibria}
\end{itemize}
\end{defbox}

\begin{exbox}[Battle of the Sexes]
Players want to coordinate but have \textbf{different preferences}:
\begin{center}
\renewcommand{\arraystretch}{1.2}
\begin{tabular}{@{} c | c c @{}}
 & Theater & Football \\
\midrule
Theater & $2, 1$ & $0, 0$ \\
Football & $0, 0$ & $1, 2$ \\
\end{tabular}
\end{center}

\textbf{Pure NE:} (Theater, Theater) and (Football, Football)
\end{exbox}

\begin{exbox}[Stag Hunt]
\begin{center}
\renewcommand{\arraystretch}{1.2}
\begin{tabular}{@{} c | c c @{}}
 & Stag & Hare \\
\midrule
Stag & $4, 4$ & $1, 3$ \\
Hare & $3, 1$ & $3, 3$ \\
\end{tabular}
\end{center}

\textbf{Story:} Two hunters can cooperate for a stag (high reward) or safely catch a hare alone.

\textbf{Pure NE:} (Stag, Stag) and (Hare, Hare)---plus mixed equilibria (see below).
\end{exbox}

\bigseparator

% --- Mixed Strategies ---
\subsubsection{Mixed Strategies}

\begin{defbox}[Mixed Strategy]
A \textbf{mixed strategy} is a probability distribution over pure strategies.

\begin{itemize}
  \item $\Delta_i = \Delta(S_i)$ = set of probability distributions over $S_i$
  \item A pure strategy is a mixed strategy with probability 1 on one action
  \item Even with finite $S_i$, there are infinitely many mixed strategies
\end{itemize}
\end{defbox}

\begin{defbox}[Expected Utility]
When mixed strategy profile $\tilde{s} \in \Delta$ is played:
\[
U_i(\tilde{s}) = \sum_{s \in S} U_i(s) \cdot \Pr(s | \tilde{s})
\]

where:
\[
\Pr(s | \tilde{s}) = \prod_{j \in N} \tilde{s}_j(s_j)
\]
(product of individual probabilities, assuming independence)
\end{defbox}

\begin{thmbox}[Nash's Existence Theorem (1951)]
Every finite game (finite players, finite action sets) has at least one \textbf{mixed strategy Nash equilibrium}.
\end{thmbox}

\separator

% --- Computing Mixed NE ---
\subsubsection{Computing Mixed Strategy Nash Equilibria}

\begin{thmbox}[Indifference Principle]
In a mixed strategy NE where player $i$ randomises over actions $A \subseteq S_i$:

\textbf{Player $i$ must be indifferent among all actions in $A$:}
\[
\forall a, a' \in A: \quad U_i(a, \tilde{s}_{-i}) = U_i(a', \tilde{s}_{-i})
\]

\textbf{Key insight:} The \textit{other} player's mixing probabilities make you indifferent!
\end{thmbox}

\begin{exbox}[Example: Matching Pennies]
\begin{center}
\renewcommand{\arraystretch}{1.2}
\begin{tabular}{@{} c | c c @{}}
 & H & T \\
\midrule
H & $1, -1$ & $-1, 1$ \\
T & $-1, 1$ & $1, -1$ \\
\end{tabular}
\end{center}

\textbf{No pure strategy NE} (verify: every cell has a profitable deviation).

\textbf{Mixed NE:} Let Player 1 play H with probability $p$, Player 2 play H with probability $q$.

\textbf{Player 2's indifference:}
\[
U_2(\text{H}) = U_2(\text{T}) \implies -p + (1-p) = p - (1-p) \implies p = \frac{1}{2}
\]

\textbf{By symmetry:} $q = \frac{1}{2}$

\textbf{Mixed NE:} $\left(\frac{1}{2}, \frac{1}{2}\right), \left(\frac{1}{2}, \frac{1}{2}\right)$ with expected utility $0$ for both.
\end{exbox}

\begin{exbox}[Example: Stag Hunt Mixed NE]
\begin{center}
\renewcommand{\arraystretch}{1.2}
\begin{tabular}{@{} c | c c @{}}
 & Stag & Hare \\
\midrule
Stag & $4, 4$ & $1, 3$ \\
Hare & $3, 1$ & $3, 3$ \\
\end{tabular}
\end{center}

Let $p = \Pr(\text{Player 1 plays Stag})$, $q = \Pr(\text{Player 2 plays Stag})$.

\textbf{Player 2's indifference:}
\[
U_2(\text{Stag}) = U_2(\text{Hare}) \implies 4p + 1(1-p) = 3p + 3(1-p)
\]
\[
4p + 1 - p = 3p + 3 - 3p \implies 3p + 1 = 3 \implies p = \frac{2}{3}
\]

\textbf{By symmetry:} $q = \frac{2}{3}$

\textbf{Three NE total:} Two pure + one mixed at $\left(\frac{2}{3}, \frac{2}{3}\right)$.
\end{exbox}

\begin{tipbox}[Exam Tip: Mixed NE Computation]
\begin{enumerate}
  \item Set up probability variables ($p$, $q$, etc.)
  \item Write expected utility equations for each action
  \item Set utilities equal (indifference condition)
  \item Solve for the \textit{other} player's probability
  \item Verify the solution makes sense (probabilities in $[0,1]$)
\end{enumerate}

\textbf{Common error:} Solving for your own probability instead of the opponent's!
\end{tipbox}

\bigseparator

% --- Maximin and Minimax ---
\subsubsection{Maximin and Minimax Strategies}

\begin{defbox}[Maximin (Safety Level) Strategy]
The \textbf{maximin strategy} for player $i$ maximises their worst-case payoff:
\[
s_i^{LS} \in \arg\max_{s_i} \min_{s_{-i}} U_i(s_i, s_{-i})
\]

The \textbf{maximin value} (safety level) is:
\[
U_i^{LS} = \max_{s_i} \min_{s_{-i}} U_i(s_i, s_{-i})
\]

\textbf{Intuition:} If opponents play to hurt you as much as possible, what's your best guarantee?
\end{defbox}

\begin{defbox}[Minimax Strategy (2-Player Games)]
The \textbf{minimax strategy} for player $i$ against player $-i$ minimises $-i$'s maximum payoff:
\[
s_i^{MM} \in \arg\min_{s_i} \max_{s_{-i}} U_{-i}(s_i, s_{-i})
\]

\textbf{Intuition:} How can you punish the opponent most, ignoring your own payoff?
\end{defbox}

\begin{exbox}[Maximin in Prisoner's Dilemma]
\begin{center}
\renewcommand{\arraystretch}{1.2}
\begin{tabular}{@{} c | c c @{}}
 & C & D \\
\midrule
C & $-1, -1$ & $-10, 0$ \\
D & $0, -10$ & $-5, -5$ \\
\end{tabular}
\end{center}

\textbf{Row player's analysis:}
\begin{itemize}
  \item If commit to C: opponent plays D (minimises row), payoff = $-10$
  \item If commit to D: opponent plays D (minimises row), payoff = $-5$
\end{itemize}

\textbf{Maximin:} Choose D (guarantees $-5 > -10$)

\textbf{Defect is both the dominant strategy AND the maximin strategy.}
\end{exbox}

\separator

% --- Zero-Sum Games ---
\subsubsection{Zero-Sum Games and the Minimax Theorem}

\begin{defbox}[Constant-Sum and Zero-Sum Games]
A two-player game is \textbf{constant-sum} if $\exists c$ such that:
\[
\forall s \in S_1 \times S_2: \quad U_1(s) + U_2(s) = c
\]

For $c = 0$: \textbf{zero-sum game} (pure conflict---one player's gain is the other's loss).
\end{defbox}

\begin{thmbox}[Minimax Theorem (von Neumann, 1928)]
For any finite, two-player, zero-sum game:
\[
\max_{s_1} \min_{s_2} U_1(s_1, s_2) = \min_{s_2} \max_{s_1} U_1(s_1, s_2)
\]

\textbf{Consequences:}
\begin{enumerate}
  \item Each player's \textbf{maximin value = minimax value}
  \item The common value is called the \textbf{value of the game}
  \item \textbf{Maximin strategies = minimax strategies = NE strategies}
  \item All Nash equilibria have the \textbf{same payoff vector}
\end{enumerate}
\end{thmbox}

\begin{exbox}[Matching Pennies: Minimax Analysis]
\begin{center}
\renewcommand{\arraystretch}{1.2}
\begin{tabular}{@{} c | c c @{}}
 & H & T \\
\midrule
H & $1, -1$ & $-1, 1$ \\
T & $-1, 1$ & $1, -1$ \\
\end{tabular}
\end{center}

\textbf{This is a zero-sum game} ($U_1 + U_2 = 0$ always).

\textbf{Value of the game:} 0

\textbf{By Minimax theorem:} The unique NE $\left(\frac{1}{2}, \frac{1}{2}\right)$ is also the maximin = minimax strategy for each player.
\end{exbox}

\begin{tipbox}[Computing Maximin/Minimax in General-Sum Games]
For a general-sum 2-player game:
\begin{enumerate}
  \item Player $i$'s maximin depends \textit{only} on $i$'s utilities
  \item Create a zero-sum game: keep $i$'s utilities, negate them for $-i$
  \item Solve the zero-sum game for NE
  \item NE strategies = maximin strategies
\end{enumerate}
\end{tipbox}

\bigseparator

% --- Other Solution Concepts ---
\subsubsection{Other Solution Concepts}

\begin{defbox}[Strong Equilibrium]
A profile $s$ is a \textbf{strong equilibrium} if no coalition $C \subseteq N$ has a joint improving move:
\[
\nexists C \subseteq N, s'_C \in S_C: \quad \forall i \in C: U_i(s'_C, s_{-C}) > U_i(s)
\]

\textbf{Key facts:}
\begin{itemize}
  \item Every strong equilibrium is a Nash equilibrium (consider singleton coalitions)
  \item Much more rare than NE
  \item Prisoner's Dilemma has \textbf{no} strong equilibrium
  \item Coordination game: both pure NE are strong equilibria
\end{itemize}
\end{defbox}

\begin{defbox}[Correlated Equilibrium]
Players condition actions on a shared random signal from a trusted mediator.

\begin{itemize}
  \item Every mixed NE induces a correlated equilibrium
  \item Set of correlated equilibria $\supseteq$ set of NE
  \item \textbf{Computationally easier:} Can be found via linear programming
\end{itemize}
\end{defbox}

\begin{defbox}[$\varepsilon$-Nash Equilibrium]
A profile where no player can gain more than $\varepsilon > 0$ by deviating:
\[
\forall i, \forall s'_i: \quad U_i(s_i, s_{-i}) \geq U_i(s'_i, s_{-i}) - \varepsilon
\]

Standard NE is the special case $\varepsilon = 0$.
\end{defbox}

\begin{defbox}[Rationalizable Strategies]
A strategy is \textbf{rationalizable} if it is a best response to some belief about opponents' strategies (consistent with common knowledge of rationality).

\begin{itemize}
  \item Equivalent to surviving iterated elimination of never-best-responses
  \item Every NE strategy is rationalizable
  \item Set of rationalizable strategies $\supseteq$ NE strategies
\end{itemize}
\end{defbox}

\bigseparator

% --- Classic Games Summary ---
\subsubsection{Summary: Classic Games}

\begin{center}
\renewcommand{\arraystretch}{1.4}
\begin{tabular}{@{} l c c c p{4cm} @{}}
\toprule
\textbf{Game} & \textbf{Pure NE} & \textbf{Mixed NE} & \textbf{Strong NE} & \textbf{Key Feature} \\
\midrule
Prisoner's Dilemma & 1 (D,D) & Same & None & Dominant strategy, Pareto-inefficient \\
Coordination & 2 & +1 & 2 & Multiple equilibria \\
Matching Pennies & 0 & 1 & 0 & Zero-sum, no pure NE \\
Battle of Sexes & 2 & +1 & 2 & Asymmetric preferences \\
Stag Hunt & 2 & +1 & 2 & Risk vs reward trade-off \\
\bottomrule
\end{tabular}
\end{center}

\bigseparator

% --- Worked Examples ---
\subsubsection{Worked Examples from Exercises}

\begin{exbox}[Exercise: Dividing Money]
\textbf{Setup:} Two players divide £10. Each names $x_i \in \{0, 1, \ldots, 10\}$.
\begin{itemize}
  \item If $x_1 + x_2 \leq 10$: each gets their claim (remainder destroyed)
  \item If $x_1 + x_2 > 10$ and $x_1 \neq x_2$: lower claimant gets $x_i$, other gets rest
  \item If $x_1 + x_2 > 10$ and $x_1 = x_2$: each gets £5
\end{itemize}

\textbf{Best responses:}
\begin{center}
\renewcommand{\arraystretch}{1.2}
\begin{tabular}{@{} c l @{}}
\toprule
\textbf{Other's claim} & \textbf{Best responses} \\
\midrule
0 & $\{10\}$ \\
1 & $\{9, 10\}$ \\
2 & $\{8, 9, 10\}$ \\
$\vdots$ & $\vdots$ \\
5 & $\{5, 6, 7, 8, 9, 10\}$ \\
6 & $\{5, 6\}$ \\
7 & $\{6\}$ \\
8 & $\{7\}$ \\
9 & $\{8\}$ \\
10 & $\{9\}$ \\
\bottomrule
\end{tabular}
\end{center}

\textbf{Nash equilibria:} $(5,5)$, $(5,6)$, $(6,5)$, $(6,6)$ --- any pair where both are best responses to each other.
\end{exbox}

\begin{exbox}[Exercise: Modified Stag Hunt]
\textbf{Modified payoffs:}
\begin{center}
\renewcommand{\arraystretch}{1.2}
\begin{tabular}{@{} c | c c @{}}
 & Stag & Hare \\
\midrule
Stag & $4, 4$ & $1, 3$ \\
Hare & $3, 2$ & $3, 3$ \\
\end{tabular}
\end{center}

\textbf{Whose mixing probability changes?}

Player 2's utility changed ($1 \to 2$ in Hare-Stag cell).

\textbf{Effect:} Player 2's utility is used to determine Player 1's indifference condition.

Player 1's mixing probability changes: $p = \frac{1}{2}$ (was $\frac{2}{3}$).

Player 2's mixing probability stays: $q = \frac{2}{3}$ (Player 1's utilities unchanged).
\end{exbox}

\bigseparator

% --- Other Game Types ---
\subsubsection{Other Game Types (Preview)}

\begin{center}
\renewcommand{\arraystretch}{1.3}
\begin{tabular}{@{} l p{10cm} @{}}
\toprule
\textbf{Game Type} & \textbf{Description} \\
\midrule
Extensive form & Game tree; players move in turns (sequential) \\
Repeated games & Same game played multiple times; cumulative payoffs \\
Cooperative (coalitional) & Players form coalitions; binding contracts; characteristic function over coalitions \\
\bottomrule
\end{tabular}
\end{center}

These will be covered in subsequent weeks.
